% ECONOMICS_PROBLEM: {"id": "L13-358", "title": "Platform competition, tipping and interoperability", "jel_code": "L13", "source_id": "OP-0073"}
% !TeX program = xelatex
\documentclass[11pt,a4paper]{article}
\usepackage[margin=23mm]{geometry}
\usepackage{fontspec}
\setmainfont{Latin Modern Roman}
\usepackage{amsmath,amssymb,mathtools}
\usepackage{xcolor,enumitem,booktabs,longtable,array,xurl,hyperref}
\definecolor{muted}{HTML}{53636C}
\hypersetup{colorlinks=true,urlcolor=blue,linkcolor=blue}
\setlength{\parindent}{0pt}
\setlength{\parskip}{5pt}
\newcommand{\R}{\mathbb R}
\newcommand{\E}{\mathbb E}
\newcommand{\N}{\mathbb N}
\newcommand{\ind}{\mathbf 1}
\newcommand{\Var}{\operatorname{Var}}
\newcommand{\Cov}{\operatorname{Cov}}
\newcommand{\supp}{\operatorname{supp}}
\DeclareMathOperator*{\argmin}{arg\,min}
\DeclareMathOperator*{\argmax}{arg\,max}
\newcommand{\entrymeta}[3]{{\sffamily\small\color{muted}\textbf{\texttt{#1}}\quad|\quad #2\par #3\par}}
\newcommand{\Problemref}[1]{\texttt{#1}}
\newcommand{\JELref}[2]{\texttt{#2} (JEL \texttt{#1})}
\begin{document}
\section*{Platform competition, tipping and interoperability}
\textbf{Problem ID:} \texttt{L13-358}\quad \textbf{JEL:} \texttt{L13}\par
% BEGIN ECONOMICS STATEMENT
\label{problem:OP-0073}
\label{atlas:prob:I3}

\noindent\textbf{Original identifier:} I3 (atlas item 73). \textbf{Classification:} Platform equilibrium and interoperability research program. \textbf{Original tractability label:} Hard (editorial, not a complexity result).

\subsection*{Definitions and assumptions}

Let platforms $k=1,\ldots,K$ serve user sides $a\in\{1,2\}$. Platform $k$ chooses side-specific fees $p_k^a$ and quality $q_k$. User utility depends on these choices, own type, and the accessible population on the opposite side; an interoperability matrix $\Gamma$ specifies cross-platform accessibility. Explicitly state multihoming costs, congestion, compatibility costs, data-sharing rules, and entry costs. A platform and user equilibrium jointly solves participation and platform optimization under these primitives. Let $\mathcal E(\Gamma)$ include possible multiple equilibria and $\sigma$ select one. For this welfare aggregation, require quasilinear utilities in a common monetary numeraire and include platform owners with the same unit weight as users. Welfare is then total gross user utility minus real production and compatibility costs, with fee transfers cancelling; externalities must be explicitly added. Other utility or distributional-weight specifications require their own explicit evaluator.

\subsection*{Formal research question}

For defined interoperability regimes $\Gamma_0,\Gamma_1$, characterize
\[
\Delta W^\sigma=W(e_\sigma(\Gamma_1),\Gamma_1)-W(e_\sigma(\Gamma_0),\Gamma_0),
\]
and the associated entry, quality, and concentration changes. Define inefficient dominance as an equilibrium allocation with lower welfare than an attainable allocation under the same resource and technological constraints. Determine conditions under which greater interoperability improves welfare for every admissible selection rule, and identify or bound these conditions from transaction, participation, and quality data. For dynamic tipping, specify an adjustment process and assess its attraction basins rather than inferring dynamics from static equilibrium multiplicity.

\subsection*{Interpretation and scope}

A large platform can be efficient because it realizes network benefits, or inefficient because market power and entry barriers dominate those gains. Market share alone does not decide between these cases. Interoperability changes both participation externalities and incentives to invest in quality, and may impose engineering, privacy, or security costs that the model must explicitly value. Fee changes across sides cannot be summarized by one price without accounting for which side generates network benefits. Static two-sided pricing results do not automatically apply to dynamic entry or product innovation. Identification requires separating users' direct taste for a platform from endogenous network size, for example through credible participation or compatibility variation. Regulatory implementations may bundle several changes; an empirical design must identify the whole bundle or justify exclusions isolating interoperability. A successful extension could characterize a welfare threshold in estimable primitives and validate its predicted changes after a specified mandate. The conclusion remains conditional on the feasible technology and equilibrium selection.

\subsection*{Source audit and status}

[S1]--[S3] verify foundational two-sided and multisided pricing analyses. They establish mechanisms and model results rather than a universal rule for interoperability. The formal intervention above introduces compatibility technology, welfare feasibility, and equilibrium selection explicitly; these are editorial additions. Existing answers in specific platform models mean the general wording is a broad research program, not one named conjecture whose present-day openness follows from the cited publication dates.

\subsection*{References}

\begin{enumerate}[label={[S\arabic*]},leftmargin=*]

\item Jean-Charles Rochet and Jean Tirole (2003). \emph{Platform Competition in Two-Sided Markets}. Journal of the European Economic Association 1(4), 990--1029. \url{https://doi.org/10.1162/154247603322493212}\par \textit{Source role:} Platform price-structure model. \textit{Locator:} Abstract and article metadata.

\item Mark Armstrong (2006). \emph{Competition in Two-Sided Markets}. RAND Journal of Economics 37(3), 668--691. \url{https://doi.org/10.1111/j.1756-2171.2006.tb00037.x}\par \textit{Source role:} Platform competition and multihoming model. \textit{Locator:} Abstract and article metadata.

\item E. Glen Weyl (2010). \emph{A Price Theory of Multi-sided Platforms}. American Economic Review 100(4), 1642--1672. \url{https://www.aeaweb.org/articles?id=10.1257/aer.100.4.1642}\par \textit{Source role:} Platform pricing and welfare framework. \textit{Locator:} Abstract and article metadata.

\end{enumerate}

\subsection*{Common conventions: Source mathematical conventions}

Unless an entry states otherwise, agent and firm sets are finite. State and action spaces are measurable, strategies and outcome maps are measurable, probability laws are specified, and expectations exist for the targets under discussion. Infinite-horizon discount factors lie in $(0,1)$ and discounted objectives are finite. Norms, tolerances and complexity input sizes must be specified for computational comparisons. Constants or primitives that appear locally are inputs, not empirical facts asserted by this catalogue.

\textbf{Identification.} A statistical model $m\in\mathcal M$ implies an observable law $P_m$ and target $\theta(m)$. At law $P$, its identified set is
\[
\Theta(P;\mathcal M)=\{\theta(m):m\in\mathcal M,\ P_m=P\}.
\]
Point identification holds when this set is a singleton, or equivalently when $P_{m_1}=P_{m_2}$ implies $\theta(m_1)=\theta(m_2)$ for all compatible models. Sharp bounds mean bounds attained, or approached when the boundary is not attained, by compatible models. Writing this set is a definition, not a solution: the substantive task is to establish informative restrictions and derive or estimate the set. An unrestricted-confounding model may give only trivial bounds.

\textbf{Inference.} Identification, estimability and valid uncertainty quantification are separate questions. Fix $\alpha\in(0,1)$. A confidence procedure $C_n$ over a declared law class $\mathcal P$ has asymptotic uniform coverage at level $1-\alpha$ when
\[
\liminf_{n\to\infty}\inf_{P\in\mathcal P}\Pr_P\{\theta(P)\in C_n\}\geq1-\alpha.
\]
For set-identified targets the coverage event must be defined separately, for example coverage of the full identified set. If an entry uses identification-indexed models instead of uniquely indexed laws, the same coverage convention is applied over those models. Pointwise limits do not establish uniform validity, and asymptotic validity does not guarantee finite-sample precision. Sample sizes, dependence structures and weak-identification sequences are part of each application.

\textbf{Counterfactuals.} $Y_i(a)$ is a potential outcome under a specified intervention, including its timing, implementation and versions. It is not $Y_i$ conditional on voluntarily choosing $a$. A full intervention vector permits interference; shorthand scalar treatment notation does not silently impose no interference. Assignment, selection, attrition, anticipation and measurement must be specified. A claim about an instrument's local effect is not automatically a population policy effect.

\textbf{Economic equilibrium.} A counterfactual model includes best responses, constraints, feasibility, market clearing, state transitions and expectations consistent with the stated information. When several equilibria exist, either a declared selector is an input or the target is set-valued. Comparisons across policy regimes must use the stated selection convention. Reduced-form relationships are not presumed invariant after an equilibrium-changing intervention.

\textbf{Optimization and welfare.} Optimization is over the stated feasible set. If attainment is not established, $\sup$ or $\inf$ and $\varepsilon$-optimal choices replace an asserted maximizer or minimizer. For example, with value $V=\sup_{a\in\mathcal A}W(a)<\infty$, an $\varepsilon$-optimal set is $\{a\in\mathcal A:W(a)\geq V-\varepsilon\}$ for $\varepsilon>0$. Benefits, costs, risk and health or environmental outcomes can be aggregated only using declared units and conversion weights. Interpersonal utility weights, the treatment of fiscal transfers and foreign profits, discounting, and the behavioral welfare criterion are explicit inputs. A comparative-static derivative additionally requires existence and appropriate differentiability of the selected counterfactual.

\textbf{Sources.} Local labels [S1], [S2], and so forth restart in every question. Locators identify the relevant abstract, model, section or result at the level actually checked; a bibliographic or abstract-level verification is not represented as inspection of an entire proof. Working papers are distinguished from later publications and changing versions. Source summaries are paraphrased. All URL checks and editorial formulations refer to the audit date stated above.
% END ECONOMICS STATEMENT
\end{document}
